
\section{The HiPPO Framework: High-order Polynomial Projection Operators}
\label{sec:framework}

We motivate the problem of online function approximation with projections
as an approach to learning memory representations
(Section~\ref{subsec:hippo-setup}).
Section~\ref{subsec:hippo-framework} describes the general
HiPPO framework to derive memory updates, including a precise definition of the technical problem 
we introduce, and an overview of our approach to solving it.
Section~\ref{subsec:high_order_projection} instantiates the framework to recover
the LMU and yield new memory updates (e.g.\ HiPPO-LagT), demonstrating the generality
of the HiPPO framework.
Section~\ref{sec:discretization} discusses how to convert the main continuous-time results into practical discrete versions.
Finally in Section~\ref{subsec:low_order_projection} we show how gating
in RNNs is an instance of HiPPO memory.


\subsection{HiPPO Problem Setup}
\label{subsec:hippo-setup}





Given an input function $f(t) \in \mathbb{R}$ on $t \ge 0$, many problems require operating on the cumulative \emph{history} $f_{\le t} := f(x) \mid_{x \le t}$ at every time $t \ge 0$,
in order to understand the inputs seen so far and make future predictions.
Since the space of functions is intractably large, the history cannot be perfectly memorized and must be compressed; we propose the general approach of projecting it onto a subspace of bounded dimension.
Thus, our goal is to maintain (online) this compressed representation of the history.
In order to specify this problem fully, we require two ingredients: a way to quantify the approximation, and a suitable subspace.


\paragraph{Function Approximation with respect to a Measure.}
Assessing the quality of an approximation requires defining a distance in function space.
Any probability measure $\mu$ on $[0, \infty)$ equips the space of square integrable functions with inner product
$
  \langle f, g \rangle_\mu = \int_0^\infty f(x) g(x) \d \mu(x),
$
inducing a Hilbert space structure $\mathcal{H}_\mu$ and corresponding norm $\| f \|_{L_2(\mu)} = \langle f, f \rangle_\mu^{1/2}$.

\paragraph{Polynomial Basis Expansion.}
Any $N$-dimensional subspace $\mathcal{G}$ of this function space is a suitable candidate for the approximation.
The parameter $N$ corresponds to the order of the approximation, or the size of the compression;
the projected history can be represented by the $N$ coefficients of its expansion in any basis of $\mathcal{G}$.
For the remainder of this paper, we use the polynomials as a natural basis, so that $\mathcal{G}$ is the set of polynomials of degree less than $N$.
We note that the polynomial basis is very general; for example, the Fourier basis $\sin(nx), \cos(nx)$ can be seen as polynomials on the unit circle $(e^{2\pi i x})^n$ (cf.\ \cref{sec:derivation-fourier}).
In \cref{sec:hippo-framework-details}, we additionally formalize a more
general framework that allows different bases other than polynomials by tilting the measure with another function.

\paragraph{Online Approximation.}
Since we care about approximating $f_{\le t}$ for every time $t$, we also let the measure vary through time.
For every $t$, let $\mu^{(t)}$ be a measure supported on $(-\infty, t]$ (since $f_{\le t}$ is only defined up to time $t$).
Overall, we seek some $g^{(t)}\in\mathcal{G}$ that minimizes $\|f_{\le t} - g^{(t)}\|_{L_2(\mu^{(t)})}$.
Intuitively, the measure $\mu$ controls the importance of various parts of the input domain,
and the basis defines the allowable approximations.
The challenge is how to solve the optimization problem in closed form given $\mu^{(t)}$,
and how these coefficients can be maintained online as $t \to \infty$.




\subsection{General HiPPO framework}
\label{subsec:hippo-framework}

We provide a brief overview of the main ideas behind solving this problem, which provides a surprisingly simple and general strategy for many measure families $\mu^{(t)}$.
This framework builds upon a rich history of the well-studied \emph{orthogonal polynomials} and related transforms in the signal processing literature.
Our formal abstraction (\cref{def:hippo}) departs from prior work on sliding transforms in several ways, which we discuss in detail in \cref{sec:rw-ops}.
For example, our concept of the time-varying measure allows choosing $\mu^{(t)}$ more appropriately, which will lead to solutions with qualitatively different behavior.
\cref{sec:hippo-framework-details} contains the full details and formalisms of our framework.

\paragraph{Calculating the projection through continuous dynamics.}
As mentioned, the approximated function can be represented by the $N$ coefficients of its expansion in any basis;
the first key step is to choose a suitable basis $\{g_n\}_{n < N}$ of $\mathcal{G}$.
Leveraging classic techniques from approximation theory, a natural basis is the set of orthogonal polynomials for the measure $\mu^{(t)}$,
which forms an orthogonal basis of the subspace.
Then the coefficients of the optimal basis expansion are simply $c^{(t)}_n := \langle f_{\le t}, g_n \rangle_{\mu^{(t)}}$.

The second key idea is to differentiate this projection in $t$, where
differentiating through the integral (from the inner product
$\langle f_{\leq t},  g_n \rangle_{\mu^{(t)}}$) will often lead to a self-similar relation
allowing $\frac{d}{dt} c_n(t)$ to be expressed in terms of $(c_k(t))_{k\in [N]}$ and $f(t)$.
Thus the coefficients $c(t) \in \R^N$ should evolve as an ODE, with dynamics determined by $f(t)$.

\paragraph{The HiPPO abstraction: online function approximation.}
\begin{definition}%
    \label{def:hippo}
    Given a time-varying measure family $\mu^{(t)}$ supported on $(-\infty, t]$, an
    $N$-dimensional subspace $\mathcal{G}$ of polynomials, and a continuous
    function $f \colon \mathbb{R}_{\geq 0} \to \mathbb{R}$,
    HiPPO defines a \emph{projection} operator $\proj_t$ and a \emph{coefficient extraction} operator $\coef_t$ at every time $t$, with the following properties:
    \begin{enumerate}[label=(\arabic*), topsep=0em, partopsep=0.0em, itemsep=0em, parsep=0.1em, leftmargin=2em]
      \item $\proj_t$ takes the function $f$ restricted up to time $t$,
      $f_{\leq t} \defeq f(x) \mid_{x \leq t}$, and maps it to a
      polynomial $g^{(t)} \in \mathcal{G}$, that minimizes the approximation error $\|f_{\leq t} - g^{(t)}\|_{L_2(\mu^{(t)})}$.
        \item $\coef_t: \mathcal{G} \to \mathbb{R}^N$ maps the polynomial $g^{(t)}$
        to the coefficients $c(t) \in \mathbb{R}^N$ of the basis of orthogonal
        polynomials defined with respect to the measure $\mu^{(t)}$.
    \end{enumerate}
    The composition $\coef \circ \proj$ is called $\hippo$, which is an operator
    mapping a function $f: \R_{\geq 0} \to \R$ to the optimal projection coefficients
    $c: \R_{\geq 0} \to \R^N$, i.e.\ $(\hippo(f))(t) = \coef_t(\proj_t(f))$.
\end{definition}

For each $t$, the problem of optimal projection
$\proj_t(f)$ is well-defined by the above inner products,
but this is intractable to compute naively.
Our derivations (\cref{sec:derivations}) will show that
the coefficient function $c(t) = \coef_t(\proj_t(f))$ has the form of an ODE
satisfying $\frac{d}{dt} c(t) = A(t) c(t) + B(t) f(t)$ for some $A(t) \in \R^{N \times N}$, $B(t) \in \R^{N \times 1}$.
Thus our results show how to tractably obtain $c^{(t)}$ \emph{online} by solving an
ODE, or more concretely by running a discrete recurrence.
When discretized, HiPPO takes in a sequence of real values and produces a
sequence of $N$-dimensional vectors.






Figure~\ref{fig:framework} illustrates the overall framework when we use uniform measures.
Next, we give our main results showing $\hippo$ for several concrete instantiations of the framework.

\begin{figure}
  \centering
  \includegraphics[width=\linewidth]{figures/hippo.pdf}
  \caption{
    \textbf{Illustration of the HiPPO framework.}
    (1) For any function $f$, (2) at every time $t$ there is an optimal projection $g^{(t)}$ of $f$ onto the space of polynomials, with respect to a measure $\mu^{(t)}$ weighing the past.
    (3) For an appropriately chosen basis, the corresponding coefficients $c(t)\in\R^N$ representing a compression of the history of $f$ satisfy linear dynamics.
    (4) Discretizing the dynamics yields an efficient closed-form recurrence for online compression of time series $(f_k)_{k\in\N}$.
  }
  \label{fig:framework}
\end{figure}


\subsection{High Order Projection: Measure Families and HiPPO ODEs}
\label{subsec:high_order_projection}

Our main theoretical results are instantiations of HiPPO for various measure families $\mu^{(t)}$.
We provide two examples of natural sliding window measures and the corresponding projection operators.
The unified perspective on memory mechanisms allows us to derive these closed-form solutions with the same strategy, provided in
Appendices \ref{sec:derivation-legt},\ref{sec:derivation-lagt}.
The first explains the core Legendre Memory Unit (LMU)~\citep{voelker2019legendre} update in a principled way and characterizes its limitations,
while the other is novel, demonstrating the generality of the HiPPO framework.
\cref{sec:derivations} contrasts the tradeoffs of these measures
(\cref{fig:measures}), contains proofs of their derivations, and derives
additional HiPPO formulas for other bases such as Fourier (recovering the
Fourier Recurrent Unit~\citep{zhang2018learning}) and Chebyshev.





The \textbf{translated Legendre (LegT)} measures assign uniform weight to the most recent history $[t-\theta, t]$.
There is a hyperparameter $\theta$ representing the length of the sliding window,
or the length of history that is being summarized.
The \textbf{translated Laguerre (LagT)} measures instead use the exponentially decaying measure, assigning more importance to recent history.
\begin{equation*}
    \textbf{LegT}: \mu^{(t)}(x) = \frac{1}{\theta} \mathbb{I}_{[t-\theta, t]}(x)
    \qquad
    \textbf{LagT}: \mu^{(t)}(x)
    = e^{-(t-x)} \mathbb{I}_{(-\infty, t]}(x)
    =
    \begin{cases}
      e^{x-t} & \mbox{if } x \le t \\
      0 & \mbox{if } x > t
    \end{cases}
\end{equation*}





\begin{theorem}
  \label{thm:legt-lagt}
  For LegT and LagT, the $\hippo$ operators satisfying \cref{def:hippo} are given by linear time-invariant (LTI) ODEs $\ddt c(t) = - A c(t) + B f(t)$, where $A \in \R^{N \times N}, B \in \R^{N \times 1}$:

  \small
  \begin{minipage}{.60\linewidth}
    \textnormal{\textbf{LegT}:}
    \vspace{-1em}
    \begin{equation}
      \label{eq:translated-legendre-dynamics}
      A_{nk} =
      \frac{1}{\theta}
      \begin{cases}
        (-1)^{n-k} (2n+1) & \mbox{if } n \ge k \\
        2n+1 & \mbox{if } n \le k
      \end{cases},
      \quad
      B_n = \frac{1}{\theta} (2n+1) (-1)^n
    \end{equation}
  \end{minipage}%
  \hfill%
  \begin{minipage}{.35\linewidth}
    \textnormal{\textbf{LagT}:}
    \vspace{-1em}
    \begin{equation}
      \label{eq:laguerre-dynamics}
      A_{nk} =
      \begin{cases}%
        1 & \mbox{if } n \ge k \\
        0 & \mbox{if } n < k \\
      \end{cases},
      \quad
      B_n = 1
    \end{equation}
  \end{minipage}

\end{theorem}





Equation~\eqref{eq:translated-legendre-dynamics} proves the LMU update \citep[equation (1)]{voelker2019legendre}.
Additionally, our derivation (Appendix~\ref{sec:derivation-legt}) shows that outside of the projections,
there is another source of approximation.
This sliding window update rule requires access to $f(t-\theta)$, which is no longer available;
it instead assumes that the current coefficients $c(t)$ are an accurate enough model of the function $f(x)_{x \le t}$ that $f(t-\theta)$ can be recovered.





\subsection{HiPPO recurrences: from Continuous to Discrete Time with ODE Discretization}
\label{sec:discretization}

Since actual data is inherently discrete (e.g.\ sequences and time series),
we discuss how the HiPPO projection operators can be discretized using standard techniques,
so that the continuous-time HiPPO ODEs become discrete-time linear recurrences.





In the continuous case, these operators consume an input function $f(t)$ and produce an output function $c(t)$.
The discrete time case (i) consumes an input sequence $(f_k)_{k \in \N}$, %
(ii) implicitly defines a function $f(t)$ where $f(k \cdot \Delta t) = f_k$ for some step
size $\Delta t$, (iii) produces a function $c(t)$ through the ODE dynamics,
and (iv) discretizes back to an output sequence $c_k := c(k \cdot \Delta t)$.


The basic method of discretizating an ODE $\frac{d}{dt} c(t) = u(t, c(t), f(t))$ chooses a step size $\Delta t$ and performs the discrete updates
$c(t+\Delta t) = c(t) + \Delta t \cdot u(t, c(t), f(t))$.\footnote{This is known as the Euler method, used for illustration here; our experiments use the more numerically stable Bilinear and ZOH methods.
\cref{sec:discretization-full} provides a self-contained overview of our full discretization framework.}
In general, this process is sensitive to the \emph{discretization step size} hyperparameter $\Delta t$.

Finally, we note that this provides a way to seamlessly handle timestamped data, even with missing values:
the difference between timestamps indicates the (adaptive) $\Delta t$ to use in discretization~\citep{chen2018neural}.
\cref{sec:discretization-full} contains a full discussion of discretization.




\subsection{Low Order Projection: Memory Mechanisms of Gated RNNs}
\label{subsec:low_order_projection}

As a special case, we consider what happens if we do not incorporate higher-order polynomials in the projection problem.
Specifically, if $N=1$, then the discretized version of HiPPO-LagT \eqref{eq:laguerre-dynamics} becomes
$c(t + \Delta t) = c(t) + \Delta t(-Ac(t) + Bf(t)) = (1 - \Delta t)c(t) + \Delta t f(t)$, since $A = B = 1$.
If the inputs $f(t)$ can depend on the hidden state $c(t)$ and the discretization step size $\Delta t$ is chosen adaptively (as a function of input $f(t)$ and state $c(t)$),
as in RNNs, then this becomes exactly a \emph{gated} RNN.
For instance, by stacking multiple units in parallel and choosing a specific update function,
we obtain the GRU update cell as a special case.\footnote{The LSTM cell update is similar, with a parameterization known as ``tied'' gates \cite{greff2016lstm}.}
In contrast to HiPPO which uses one hidden feature and projects it onto high order polynomials,
these models use many hidden features but only project them with degree 1.
This view sheds light on these classic techniques by showing how they can be derived from first principles.



